\section{Steady-state linear response in noisy network dynamics}
% \subsection{}
We now develop a steady-state linear-response framework for noisy network dynamics under parameter perturbations. We focus on systems that relax to a stable fixed point and fluctuate around it under Gaussian white noise. For the linearized Ornstein-Uhlenbeck process introduced in Sec. II, the response identity derived below gives the exact first-order susceptibility. When it is applied back to the original nonlinear network, it should be interpreted as the corresponding fixed-point approximation, consistent with the local Gaussian steady-state form discussed in Sec. II and Appendix~A. When the parameters are perturbed as $\vec{p} \rightarrow \vec{p} + \delta \vec{p}$, the steady-state average of an observable ${\bf B}(\vec{x}; \vec{p})$ changes accordingly. Our goal is to express the response of steady-state observables in terms of quantities measured in the unperturbed steady state. We first consider how the noise-free fixed point depends on the parameters. In Eq. (\ref{deterministic_force}), the noise-free fixed point satisfies
\begin{equation}
    \vec{F}(\vec{X}; \vec{p}) = 0.
\end{equation}
When the parameters change as $\vec{p} \rightarrow \vec{p} + \delta \vec{p}$, the fixed point shifts from $\vec{X}$ to $\vec{X}+\delta \vec{X}$. Expanding
\begin{equation}
    \vec{F}(\vec{X}+\delta \vec{X}; \vec{p}+ \delta \vec{p}) = 0
\end{equation}
to linear order gives
\begin{equation}
    \nabla \vec{F}\big|_{\vec{X}; \vec{p}} \delta \vec{X} + \nabla_{\vec{p}}\vec{F}\big|_{\vec{X}; \vec{p}} \delta \vec {p} = 0.
\end{equation}
With $\mathbf{Q}\equiv\nabla\vec{F}\big|_{\vec{X}}$ and assuming that $\mathbf{Q}^{-1}$ exists, one obtains
\begin{equation}
    \frac{\partial X_i}{\partial p_\alpha} = -\left(\mathbf{Q}^{-1} \nabla_{\vec{p}} \vec{F}\big|_{\vec{X}; \vec{p}}\right)_{i\alpha}.
\end{equation}
This relation shows that the shift of the fixed point induced by a parameter perturbation is controlled by the linearized dynamics through $\mathbf{Q}$. For an observable $\mathbf{B}(\vec{x}; \vec{p})$, the steady-state linear response is defined by
\begin{equation}
    \langle \mathbf{B} \rangle_{ss, \vec{p}+\delta \vec{p}} = \langle \mathbf{B} \rangle_{ss, \vec{p}} + \boldsymbol{\chi} \delta \vec{p}.
    \label{response}
\end{equation}
The steady-state response function is then
\begin{equation}
    \boldsymbol{\chi} = \langle \nabla_{\vec{p}} \mathbf{B} \rangle_{ss, \vec{p}} + \int d\vec{x} \mathbf{B} \nabla_{\vec{p}} \psi_{ss}.
    \label{response_function}
\end{equation}
The steady state distribution can be written as $\psi_{ss}=Z^{-1}e^{-\Phi}$, and we define the conjugate force with respect to the parameters by
\begin{equation}
    \vec{\mathscr{F}} \equiv -\nabla_{\vec{p}} \Phi.
    \label{conjugate_force}
\end{equation} 
Notice that $\Phi=\Phi(\vec{p}, \vec{X}(\vec{p}))$, where the fixed point depends on the parameters. One obtains
\begin{equation}
    \partial_{p_\alpha} \psi_{ss} = \psi_{ss}(\mathscr{F}_\alpha-\langle \mathscr{F}_\alpha \rangle).
\end{equation}
Therefore,
\begin{equation}
    \boldsymbol{\chi} = \langle \nabla_{\vec{p}} \mathbf{B} \rangle_{ss, \vec{p}} + \langle \delta \mathbf{B} \delta \vec{\mathscr{F}}^\intercal \rangle_{ss, \vec{p}}.
    \label{response_function2}
\end{equation}
This is the steady-state linear-response formula used throughout this work \cite{softmatter,Seifert_2010,Mehl_2010}. It requires a differentiable steady-state distribution near a stable fixed point, but it does not require equilibrium. In the present paper, the explicit formulas derived from it are evaluated within the linearized Gaussian description, where they give exact first-order susceptibilities for the Ornstein-Uhlenbeck dynamics and leading-order approximations for the original nonlinear system near the fixed point. In the following sections, we apply this framework to perturbations of the intrinsic node parameter $r_\alpha$, the connection weight $W_{\alpha\beta}$, and the noise strength $\sigma_{\alpha\alpha}$. The detailed derivation of the conjugate force and the related response identities is summarized in Appendix~B.
